\section{Structural Emergence of Mean Bias in Transformer Representations}

The section shows coherent mean bias observed in Transformer models is not an accidental artifact but a predictable consequence of the interaction between language statistics and the specific geometry of high-dimensional operators. We argue that this bias emerges through a three-stage causal chain: frequency-weighted initialization, architectural regeneration, and residual accumulation.

\paragraph{Statistical Origin: Frequency-Weighted Embedding Expectation.}

A ``seed'' bias is introduced at the very first layer of the model.
Let $\mathcal{V}$ denote the vocabulary and let $p(v)$ be the token frequency
distribution over $\mathcal{V}$.
In natural language corpora, $p(v)$ is highly non-uniform and typically
follows a Zipfian distribution.
The expected embedding vector over the corpus is

\begin{equation}
\mu_{\text{embed}}
=
\mathbb{E}_{v \sim p}[E_v]
=
\sum_{v \in \mathcal{V}} p(v) E_v,
\label{eq:embedding_mean}
\end{equation}

where $E_v \in \mathbb{R}^H$ denotes the embedding vector of token $v$
\cite{mikolov2013efficient,mikolov2013distributed,arora2017simple}.

Because high-frequency tokens (e.g., articles, punctuation, and syntactic
markers) dominate $p(v)$ and receive disproportionately many gradient updates
during training, their embedding vectors tend to align along common predictive
directions.
In high-dimensional spaces, such coherent alignment induces a dominant
rank-one component in the embedding matrix even before subsequent Transformer
processing begins.

\paragraph{The Mechanism of Mean Regeneration.}

Even if internal activations are perfectly centered at zero,
the fundamental operators of the Transformer architecture
systematically regenerate a coherent mean.
This regeneration is not incidental but arises from
the algebraic structure of nonlinearities,
attention mechanisms, and residual integration.

\noindent \emph{Non-Odd Nonlinearities as Bias Generators:} Most modern activation functions $\phi(\cdot)$ used in
Transformer architectures—such as ReLU, GELU, and SwiGLU—
are non-odd functions.
For any symmetric zero-mean random variable $z$ satisfying
$\mathbb{E}[z] = 0$ and $z \overset{d}{=} -z$,
non-odd nonlinearities satisfy

\begin{equation}
\mathbb{E}[\phi(z)] > 0,
\label{eq:non_odd_bias}
\end{equation}

under mild regularity conditions
\cite{timkey2021rogue,puccetti2022frequency}.

For example, ReLU is defined as $\phi(z) = \max(0,z)$.
If $z$ is symmetric, then

\begin{equation}
\mathbb{E}[\max(0,z)]
=
\int_{0}^{\infty} z\, p(z)\, dz
>
0,
\end{equation}

because all negative mass is truncated.
Thus ReLU introduces a strictly positive drift even when
the input is perfectly centered.
The same phenomenon holds approximately for GELU and SwiGLU,
which suppress negative regions asymmetrically.
Only strictly odd nonlinearities such as $\tanh$
preserve zero mean in expectation.

Therefore, unless $\phi$ is explicitly odd,
mean regeneration is mathematically unavoidable.

\noindent \emph{Softmax Attention as Convex Aggregation:}
Self-attention further reinforces mean drift.
Given query-key-value projections

\begin{equation}
Q = XW_Q,\quad K = XW_K,\quad V = XW_V,
\end{equation}

the attention output is

\begin{equation}
Y_i = \sum_{j} A_{ij} V_j,
\quad
A = \mathrm{softmax}\!\left(\frac{QK^\top}{\sqrt{d}}\right).
\end{equation}

Because softmax enforces positivity and row-wise normalization,
$A_{ij} \ge 0$ and $\sum_j A_{ij} = 1$,
each output $Y_i$ is a convex combination of value vectors.

Decompose values as

\begin{equation}
V_j = c + \epsilon_j,
\end{equation}

where $c$ is a shared component and $\epsilon_j$ are deviations
with $\sum_j \epsilon_j = 0$.
Then

\begin{equation}
Y_i
=
\sum_j A_{ij}(c + \epsilon_j)
=
c + \sum_j A_{ij}\epsilon_j.
\label{eq:attention_mean_preserve}
\end{equation}

The shared component $c$ is preserved exactly.
Meanwhile, the deviations $\epsilon_j$ are averaged.
Because $A_{ij}$ depends on $QK^\top$,
which itself depends on $V$ through previous layers,
the symmetry conditions required for cancellation are violated.
Thus attention acts as a data-dependent low-pass filter,
preserving coherent components while attenuating fluctuations.

\paragraph{Residual Integration and High-Dimensional Magnification.}

Mean drift generated by nonlinearities and attention
is integrated through residual connections:

\begin{equation}
\mu_{l+1}
=
\mu_l + \Delta\mu_l,
\label{eq:residual_accumulation}
\end{equation}

where $\Delta\mu_l$ denotes the mean generated by layer $l$.

This additive coupling causes bias to accumulate approximately linearly
with depth when $\Delta\mu_l$ maintains consistent orientation.

Crucially, high dimensionality magnifies this effect.
Let $\bar{\mu}$ denote the average coordinate-wise bias.
Then the Euclidean norm of the mean vector scales as

\begin{equation}
\|\mu\|_2
\sim
\sqrt{H}\,\bar{\mu},
\label{eq:mean_scaling}
\end{equation}

where $H$ is the hidden dimension
\cite{tropp2012user,halko2011randomized}.

In a space with $H=4096$,
a seemingly small per-coordinate bias is amplified
by a factor of $\sqrt{4096}=64$ in total magnitude.
This geometric magnification explains why mean-dominated components
account for a large fraction (empirically $85\%$–$90\%$)
of extreme activation outlier energy.

Thus, mean regeneration arises from:
(i) non-odd nonlinearities,
(ii) convex aggregation in softmax attention,
and (iii) residual accumulation amplified by high dimensionality.


\paragraph{Empirical Verification.}

We measure mean magnitude using the normalized ratio

\begin{equation}
R = \frac{\|\mu\|_2}{\sqrt{\mathbb{E}\|X\|^2/N}}.
\end{equation}

Across embedding, shallow, middle, and deep layers,
and across early and late training stages:

\begin{itemize}
\item Activation $R$ ranges from $10$ to $200+$.
\item Gradient-output $R \approx 1$.
\item Mean bias strengthens over training.
\item Mean bias increasingly dominates top $0.1\%$ activation magnitudes.
\end{itemize}

Additionally, decomposing activations into
mean + spike + tail components shows that:

\begin{equation}
\text{Outlier Energy}_{\text{mean}}
\approx
85\% \text{–} 90\%.
\end{equation}

These findings confirm the theoretical predictions:
mean bias is the principal driver of forward-path anisotropy.

